\section*{Hoofdstuk \ref{ch:g12:past-climates} --- Klimaten van het verleden lezen}

\begin{solution}{exo:g12:past-climates:1}
(a) een \omterm{def:g12:past-climates:proxy}{proxy}; (b) een \omterm{def:g12:past-climates:ice-core}{ijskern}; (c) een \omterm{def:g12:past-climates:glacial-cycle}{interglaciaal}; (d) \omterm{def:g12:past-climates:pollen}{pollenanalyse};
(e) het \omterm{def:g12:past-climates:lgm}{Laatste Glaciale Maximum}.
\end{solution}

\begin{solution}{exo:g12:past-climates:2}
(a) Polair ijs: het koolstofdioxide en het methaan van de lucht van toen.
(b) Meersedimenten of veen: de begroeiing, hier bomen die warme
zomers nodig hebben. (c) Oceaansedimenten: het volume ijs op het land en de temperatuur
van de \omterm{def:g10:layered-ocean:thermocline}{diepe oceaan}. (d) Bomen: de groeiomstandigheden, jaar na jaar.
\end{solution}

\begin{solution}{exo:g12:past-climates:3}
(a) $-50 - (-120) = \qty{70}{m}$ in \num{7000} jaar: $\qty{70000}{mm}/\num{7000}
= \qty{10}{mm}$ per jaar. (b) $10/3.7 \approx 2.7$: de gemiddelde stijging van de
deglaciatie was ongeveer 2.7 keer sneller dan die van 2006--2018.
\end{solution}

\begin{solution}{exo:g12:past-climates:4}
Een eeuw is tien decennia: $\qty{2.0}{\degreeCelsius}$ per eeuw. Duizend
jaar is honderd decennia: $\qty{20}{\degreeCelsius}$ per duizend jaar.
\end{solution}

\begin{solution}{exo:g12:past-climates:5}
(a) \num{23000} jaar: de precessie van de as; \num{42000} jaar: de
helling van de as; ongeveer \num{100000} jaar: de vorm van de baan.
(b) $\num{1000000}/\num{41000} \approx 24$ cycli; $\num{1000000}/\num{100000}
= 10$ cycli.
\end{solution}

\begin{solution}{exo:g12:past-climates:6}
(a) $266.6 + 7.83 \times (-9.4) \approx 266.6 - 73.6 = \qty{193}{ppm}$: maar
\qty{3}{ppm} boven de gemeten \qty{190}{ppm}. (b) $266.6 + 7.83 \times 5.5
\approx \qty{310}{ppm}$, \qty{11}{ppm} boven de hoogste gemeten waarde: het
warmste punt ligt onder de lijn. (c) $\Delta T = (427 - 266.6)/7.83 \approx
\qty{20}{\degreeCelsius}$. De lijn beschrijft alleen hoe de twee grootheden
samen varieerden, in natuurlijke cycli, binnen het bereik van de gegevens
(\qty{-11}{\degreeCelsius} tot \qty{+6}{\degreeCelsius}, 174 tot
\qty{299}{ppm}); ze is geen wet, en een waarde ver buiten dat bereik kan
er niet uit worden afgelezen.
\end{solution}

\begin{solution}{exo:g12:past-climates:7}
(a) Kruiden, \qty{84}{\percent}: een open landschap met bijna geen bomen, koud.
(b) Eik (\qty{42}{\percent}) en hazelaar (\qty{34}{\percent}): een loofbos,
met warme zomers. (c) De berk breidt zich als eerste uit, dan de den; hazelaar en eik
komen het laatst. (d) Bomen breiden zich maar zo snel uit als hun zaden reizen en hun
jonge bomen groeien, vanaf de plaatsen waar ze de kou overleefden; er moet
ook een bodem ontstaan. Het bos kan achterblijven bij het klimaat.
\end{solution}

\begin{solution}{exo:g12:past-climates:8}
(a) \qty{300}{cm}, dat is \qty{3000}{mm}, in \num{8000} jaar: $3000/8000 =
\qty{0.375}{mm}$ per jaar. (b) $\num{6000} + (470 - 300)/300 \times \num{8000}
\approx \num{6000} + \num{4530} \approx \num{10500}$ jaar. (c) $10/0.375
\approx 27$ jaar.
\end{solution}

\begin{solution}{exo:g12:past-climates:9}
(a) $7.16 - 4.11 = 3.05$ miljoen vierkante kilometer, en $3.05/4.11 \approx
\qty{74}{\percent}$ meer land. (b) Ja: tijdens het \omterm{def:g12:past-climates:lgm}{Laatste Glaciale Maximum} stond de zee
125 tot \qty{134}{m} lager dan vandaag, veel meer dan de \qty{36}{m}
die op de kaart Borneo met Azië verbinden. (c) De kaart gebruikt de zeebodem van vandaag, die sindsdien niet
precies dezelfde is gebleven (er zijn sedimenten op afgezet, en de
grond is een beetje bewogen); het rooster is ook grof, één punt om de 0.1
graad, en het laagste zeeniveau is maar tot op ongeveer tien meter bekend.
\end{solution}

\begin{solution}{exo:g12:past-climates:10}
(a) De waarde van het \omterm{def:g12:past-climates:lgm}{Laatste Glaciale Maximum}, \qty{4.92}{\permil}, is de
hoogste. Dat is te verwachten: de grote ijskappen van toen sloegen water op dat arm was
aan zuurstof-18, waardoor de oceaan, en de schelpen die erin groeiden, rijker werden aan
zuurstof-18; het diepe water was ook kouder. (b) Zodat koude, ijsrijke tijden
laag liggen en warme tijden hoog, zoals op een temperatuurkromme.
\end{solution}

\begin{solution}{exo:g12:past-climates:11}
(a) Een correlatie toont dat de twee reeksen samen stijgen en dalen, dicht
bij een rechte; ze toont niet welke welke veroorzaakte, en ook niet dat een van beide dat
deed. (b) Een warmer klimaat kan koolstofdioxide in de lucht hebben vrijgemaakt;
of een derde oorzaak, zoals de veranderingen van de baan van de Aarde, kan
beide hebben bewogen; of elk kan het andere hebben aangeduwd. (c) De timing van de veranderingen,
welke eerst kwam, uit een nauwkeurige datering van het ijs en van zijn lucht; en
inzicht in hoe koolstofdioxide op de straling van de Aarde inwerkt
(\cref{ch:g12:climate-change}).
\end{solution}

\begin{solution}{exo:g12:past-climates:12}
(a) $\num{1300000}/\num{41000} \approx 32$ cycli tussen 2.6 en 1.3
miljoen jaar geleden; $\num{700000}/\num{98500} \approx 7$ cycli over de laatste
\num{700000} jaar. (b) \num{41000} jaar: de helling van de as; ongeveer
\num{100000} jaar: de vorm van de baan. (c) Een cyclus van \num{41000} jaar
wordt ongeveer $\num{41000}/\num{2000} \approx 20$ keer bemonsterd, genoeg om hem te tekenen;
een verandering die een eeuw duurt, valt tussen twee monsters, of binnen één, en
laat bijna geen spoor na.
\end{solution}

\begin{solution}{exo:g12:past-climates:13}
(a) $192/10 \approx 19$ en $198/10 \approx 20$: ongeveer 20 keer sneller.
(b) Tegenover \qty{4}{ppm} per eeuw: $192/4 = 48$ tot $198/4 \approx 50$
keer; tegenover \qty{42}{ppm}: $192/42 \approx 4.6$ tot $198/42 \approx 4.7$
keer. De stijging van vandaag is ongeveer 5 tot 50 keer sneller dan die van het PETM.
(c) De stijging is $2000 - 900 = \qty{1100}{ppm}$: $1100/42 \approx 26$ eeuwen,
ongeveer \num{2600} jaar bij de snelste snelheid van het PETM; $1100/195 \approx 5.6$
eeuwen, minder dan 600 jaar, bij de snelheid van vandaag.
\end{solution}

\begin{solution}{exo:g12:past-climates:14}
(a) Tien keer 1 tot \qty{1.5}{\degreeCelsius}: 10 tot \qty{15}{\degreeCelsius}
per duizend jaar, dat is 0.10 tot \qty{0.15}{\degreeCelsius} per decennium.
(b) De mondiale snelheid van vandaag, \qty{0.20}{\degreeCelsius} per decennium, is 1.3 tot 2
keer de snelste opwarming boven Groenland. (c) Eén plaats kan veel
meer veranderen dan het gemiddelde over de aardbol, waar de veranderingen van verschillende streken
elkaar deels opheffen: een mondiale snelheid wordt vergeleken met een mondiale snelheid.
\end{solution}

\begin{solution}{exo:g12:past-climates:15}
(a) \qty{14}{m} aan \qty{12}{m} per duizend jaar duurt $14/12 \approx
1.17$ duizend jaar, ongeveer \num{1170} jaar: ongeveer \num{9800} jaar geleden.
(b) Nee: de zee moest minstens \qty{36}{m} lager blijven dan vandaag om
Borneo met Azië te verbinden; zodra ze \qty{-36}{m} passeerde, liep de landbrug onder, en
was Borneo weer een eiland. (c) $36/12 = 3$ duizend jaar, ongeveer
\num{3000} jaar; in werkelijkheid vertraagde de stijging na ongeveer \num{8200} jaar
geleden, dus duurde het langer.
\end{solution}

\begin{solution}{pb:g12:past-climates:1}
\textbf{1.} Het ijs zelf legt via zijn zware isotopen de
temperatuur vast; de lucht in zijn belletjes legt het koolstofdioxide vast.

\textbf{2.} $5.5 - (-10.6) = \qty{16.1}{\degreeCelsius}$.

\textbf{3.} $299 - 174 = \qty{125}{ppm}$.

\textbf{4.} De punten liggen dicht bij een rechte: de twee reeksen stegen
en daalden samen. Hij zegt niet welke welke veroorzaakte.

\textbf{5.} Reeksen uit het verleden worden getekend met de oudste tijden aan de ene of de andere
kant; verkeerd om gelezen, zou een opwarming voor een afkoeling worden gehouden.

\textbf{6.} Een \omterm{def:g12:past-climates:glacial-cycle}{glaciaal} is een koud tijdvak binnen een ijstijd, waarin de
ijskappen ver boven hun huidige omvang uitgroeien; een \omterm{def:g12:past-climates:glacial-cycle}{interglaciaal} is een warm
tijdvak tussen twee glacialen, met ijskappen ongeveer zoals vandaag of
kleiner.

\textbf{7.} In een \omterm{def:g12:past-climates:glacial-cycle}{interglaciaal}, dat ongeveer \num{11000} jaar geleden begon.

\textbf{8.} $\num{800000}/\num{98000} \approx 8$: ongeveer acht cycli.

\textbf{9.} De vorm van de baan.

\textbf{10.} Een ritme van \num{41000} jaar, de periode van de helling van de
as.

\textbf{11.} $-0.4 - (-9.4) = \qty{9.0}{\degreeCelsius}$.

\textbf{12.} $264 - 190 = \qty{74}{ppm}$.

\textbf{13.} \qty{6}{\degreeCelsius}; de verandering in Antarctica, \qty{9.0}{\degreeCelsius},
is groter, ongeveer 1.5 keer de mondiale.

\textbf{14.} $\num{18000} - \num{11000} = \num{7000}$ jaar.

\textbf{15.} $6/7 \approx \qty{0.86}{\degreeCelsius}$ per duizend jaar.

\textbf{16.} Duizend jaar is honderd decennia: $0.86/100 \approx
\qty{0.0086}{\degreeCelsius}$ per decennium.

\textbf{17.} $0.20 \times 100 = \qty{20}{\degreeCelsius}$ per duizend jaar.

\textbf{18.} $20/0.86 \approx 23$ keer.

\textbf{19.} $20/1.5 \approx 13$ en $20/1 = 20$: 13 tot 20 keer.

\textbf{20.} De mondiale opwarming van vandaag is \textbf{ongeveer twintig keer sneller} dan
de opwarming uit de laatste ijstijd: 23 keer haar gemiddelde snelheid, 13 tot 20 keer
haar snelste.
\end{solution}
